\section{Introduction and Current State}
\label{sec:introduction}

\plantodo{S1}{Briefly state the subject, purpose, and scope of the report:
the NAS modeled as a System of Systems, centered on how operational intent
propagates from enterprise objectives to aircraft trajectory/behavior.}

My project aims to build a reference MBSE architecture for how airspace intent is managed across the NAS — spanning the organizations that define and enforce rules, the domains managing airspace resources, the systems supporting flight operations, and the onboard systems that turn intent into executable trajectory and guidance behavior.

In his 2005 paper, DeLaurentis states the following~\textcite{delaurentis2005sosTransportation}:

\begin{quote}
``PROBLEMS of significant complexity are facing decision-makers within government and industry, and the moniker of `system-of-systems' (heretofore shortened to SoS) is increasingly being applied to problems of this type. Effective system analysis for decision-support quickly becomes unmanageable in this context, with multiple, heterogeneous, distributed systems involved and couched in networks at multiple levels. The situation is exacerbated by the fact that most organizations (and their processes) remain in the `stovepipe' mode -- their people and tools are configured to study within narrow bins with little or no analysis across those bins. Current frames of reference, thought processes, analysis, and design methods are not complete for these SoS problems, exemplified in this paper by future national transportation system realizations. A holistic framework is needed that enables decision makers to discern whether related infrastructure, policy, and/or technology considerations together are good, bad (or indifferent) over time. And this need is urgent, for the SoS type almost always involves decisions that commit large amounts of money, for which ultimate failure or success carries heavy consequences over several generations, and that impact large segments of the public. The costs associated with a chosen path may still be high, but the intent is to maximize the probability that these costs result in good consequences rather than bad.''
\end{quote}

While this reference is from 2005, it is still relevant today. The NAS is a complex system of systems, and the decisions made by the various stakeholders have significant consequences for the safety, efficiency, and sustainability of the air transportation system.

Airline companies invest significant time and money in constructing operational plans that maximize revenue; however, on the day of operation, unpredictable events -- bad weather, aircraft malfunctions, crew absenteeism -- can substantially alter that plan and the revenue objectives behind it~\textcite{castro2013aoccMasThesis}. To manage this, airlines maintain an Airline Operations Control Center (AOCC), a department of experts specialized in disruption management across aircraft, crew, and passenger recovery. The AOCC monitors operations and makes real-time decisions -- re-routing flights, rescheduling crews, rebooking passengers -- intended to minimize the impact of a disruption on the airline's schedule and revenue objectives.

The literature does not converge on a single, standardized definition of the "cost" the AOCC is minimizing, and the values it does quantify do not always hold up under scrutiny. Delay cost is commonly modeled in the optimization literature as a linear function of delay minutes, but \textcite{hu2024disruptionOptReview} note this assumption is inconsistent with airline practice, where delay cost is in fact a nonlinear function of delay time -- and that, faced with this complexity, airlines tend to prioritize the speed of a recovery decision over the mathematical optimality of the resulting solution. This preference for tractable, fast decisions over globally optimal ones shows up structurally as well as behaviorally: a systematic review of the Aircraft Recovery Problem finds that studies -- and, by extension, the AOCC processes they model -- rarely minimize delay, cancellations, and aircraft swaps simultaneously, instead treating aircraft, crew, and passenger recovery as separable sub-problems solved sequentially and handed off in turn~\textcite{santana2023arpReview}. Each stage can produce a solution that is locally optimal yet globally suboptimal, or even infeasible, once the next stage's constraints are applied -- a concrete, literature-documented instance of the myopic-optimization pattern this capstone investigates.

That pattern's costs and benefits are visible even in work that improves on the status quo. \textcite{kontodimou2026turnaroundBuffer}'s predictive turnaround-buffer allocation model achieves a 37.1\% reduction in propagated delay relative to a zero-buffer baseline,but it reaches that result by holding aircraft routing, fleet assignment, crew duty, maintenance, and passenger connectivity fixed as upstream constraints, deliberately excluding them from the optimization it performs. The gain is real within that bounded scope, but the framing illustrates the exact risk this capstone is concerned with: a locally-optimized decision layer can report a strong local benefit while externalizing tradeoffs onto adjacent decision layers that go unmeasured by that layer's own cost model. The same pattern appears on the human side of the AOCC: \textcite{dispatcherWorkload2025} documents that when dispatcher workload exceeds capacity, dispatchers "may resort to conservative approaches, which can have financial implications for airlines" -- a workload-driven cost that rarely appears in the delay, cancellation, and swap cost models discussed above. \textcite{hassanDisruptionReview} frames this more broadly as a persistent gap between the disruption-management optimization literature and the reality of AOCC practice.

In this paper, I explore the value tradeoffs that aviation stakeholders experience when disruptions occur and recovery solutions are employed with various levels of myopia. When a dispatcher or AOCC team member makes a decision, they are often focused on the immediate impact of that decision -- e.g., minimizing delay for a specific flight -- without visibility into the broader implications for the airline's overall operations or the NAS as a whole. Cost models used today typically capture the factors that are easiest to quantify, such as fuel consumption and crew costs, while factors that are harder to quantify -- maintenance cost, passenger satisfaction, dispatcher/controller workload, environmental impact -- are more often overlooked. This can lead to suboptimal outcomes: increased congestion in certain airspace sectors, missed opportunities for more efficient routing, or decreased overall profits driven by customer dissatisfaction. By modeling the NAS as a system of systems and analyzing the propagation of operational intent, I aim to identify strategies for improving decision-making and reducing the negative impacts of disruptions on the air transportation system.


\subsection{Overview of the Project}
\label{sec:overview}

This project exists within the industry and operating environment of U.S.
commercial air transportation -- specifically, the National Airspace System
(NAS): the combination of airspace, airports, aircraft, air navigation
facilities, and the regulatory and operational organizations that use them to
move passengers and cargo. The NAS is not one system but, as introduced above,
a System of Systems: it is jointly owned and operated by organizations with
distinct, only partially overlapping objectives -- airlines, airports, the
FAA and its Air Traffic Control (ATC) facilities, and the aircraft and crews
that execute the plans those organizations produce. No single one of these
organizations has authority over the whole; each manages its own slice of the
problem and hands intent -- a schedule, a flight plan, a clearance, a
trajectory -- to the next.

The social context surrounding this environment is significant on several
axes at once. It is safety-critical: air traffic control exists principally
to guarantee separation between aircraft, and the flight crew's foremost
authority is to deviate from any assigned plan to preserve that safety margin.
It is economically significant: a single major carrier's irregular
operations -- weather, mechanical failures, and crew unavailability -- can
exceed \$440 million per year in lost revenue, crew overtime, and passenger
accommodation costs, and a single severe-weather event (the "Blizzard of
'96") cost the U.S. airline industry an estimated \$50-100 million in one
week~\textcite{clarke1998irregular}. It depends on specialized human
operators -- dispatchers and controllers -- whose workload is itself a
scarce, unpriced resource: dispatchers facing workload beyond their capacity
"may resort to conservative approaches, which can have financial implications
for airlines"~\textcite{dispatcherWorkload2025}, and controller workload is a
hard operational ceiling on how much traffic a sector can safely
absorb. And it has environmental and community consequences -- fuel burn,
emissions, and aircraft noise -- that are borne by parties (future
generations, communities near airports) who are not party to the operational
decisions that produce them. Section~\ref{sec:stakeholders} develops this
stakeholder landscape in full; the point to establish here is that the
NAS's operating environment already contains, in miniature, the tension this
project investigates: safety, economics, human workload, and environmental
cost are managed by different organizations, on different timescales, using
different and sometimes incompatible definitions of "good."

I selected this topic for two reasons. First, the NAS is one of the most
mature, best-documented examples of a real, operating System of Systems
available to a systems-engineering capstone: the FAA and its international
counterpart, EUROCONTROL/SESAR, publish detailed public planning and
architecture documents (e.g.,
\textcite{faaNasInfrastructureRoadmaps2025,sesarju2025masterPlan}), and
decades of operations-research and disruption-management literature document
how its constituent organizations actually make decisions. That makes it
possible to build a grounded architecture rather than a purely hypothetical
one. Second, and more specifically, this project began as a narrower
trajectory-optimization problem and was deliberately broadened into an
architecture of the whole intent-propagation chain (see
\hyperref[sec:scope]{Scope} and the project's decision log, D-001) because
that framing better exercises the skills an ISD capstone is meant to
demonstrate -- structure, interfaces, responsibility, and traceability --
while still leaving room for an optimization case study that shows why those
skills matter operationally.

The potential benefit of this work is a reference architecture that makes
cross-stakeholder tradeoffs visible \emph{before} they are made, not after.
As the literature reviewed in Section~\ref{sec:litreview} shows repeatedly,
today's decision-support tools and human decision processes are organized
around one organization's cost function at a time; the costs that decision
externalizes onto another stakeholder are usually invisible to the tool or
person making it. An architecture that traces how operational intent and
information actually flow -- from an enterprise objective, through airline
and ATC planning, to the trajectory an aircraft actually flies -- gives a
systems engineer, and eventually a decision-support capability built on that
architecture, a place to see those externalities before they are absorbed as
unplanned cost by someone else.

Six questions motivate the work and give it a testable shape (see also the
project's D-004 decision record):

\begin{itemize}
  \item \emph{How should the NAS be decomposed into constituent systems and
    Systems-of-Systems?} -- the architecture is only useful if its boundaries
    reflect real authority and information-ownership lines, not an arbitrary
    catalog of aviation objects.
  \item \emph{Where do authority, responsibility, ownership, and lifecycle
    boundaries occur?} -- these are exactly the transitions (e.g., dispatcher
    to controller, controller to controller, controller to flight crew) where
    intent can be lost, delayed, or reinterpreted.
  \item \emph{How does operational intent propagate from enterprise
    objectives to physical aircraft behavior?} -- this is the project's
    organizing throughline (Strategic objective $\rightarrow$ mission plan
    $\rightarrow$ flight plan $\rightarrow$ ATC constraints $\rightarrow$
    trajectory negotiation $\rightarrow$ FMS intent $\rightarrow$ guidance
    commands $\rightarrow$ aircraft motion) and the spine the architecture is
    built around.
  \item \emph{How do stakeholder definitions of "optimal" differ?} --
    Section~\ref{sec:stakeholders} and the project's objective/cost ontology
    show that "optimal" means something different, and sometimes
    incompatible, to an airline, a controller, a passenger, and a community
    near an airport.
  \item \emph{When does subsystem optimization produce undesirable
    System-of-Systems outcomes?} -- this is the myopic-optimization pattern
    documented throughout Section~\ref{sec:litreview}, and the specific
    failure mode this project is trying to expose rather than hide.
  \item \emph{How can an MBSE architecture support analysis of those
    tradeoffs?} -- this is the project's methodological bet: that a
    traceable model, not a one-off analysis, is what lets a tradeoff
    surfaced in one scenario generalize to the next.
\end{itemize}

Taken together, these questions state the problem this project addresses:
NAS decision-making is fragmented across specialized organizational silos --
airline planning, day-of-operations control, air traffic control, airport
operations -- each of which optimizes against its own, usually well-defined
cost model, while the couplings \emph{between} those silos are rarely modeled
explicitly and are usually only discovered after a disruption forces them
into view. That matters because, as
Section~\ref{sec:litreview} documents, this is not a hypothetical risk: it is
the specific, recurring pattern the disruption-management and airline-
scheduling literatures describe when they find that sequential, decomposed
optimization is common in practice and produces solutions that are locally
optimal but globally suboptimal or infeasible once handed to the next stage.
Building a reference architecture that makes that propagation chain explicit
is the first step toward analyzing, rather than merely enduring, the
tradeoffs it creates.

\subsection{Review of Current State Design / Literature Review}
\label{sec:litreview}

As of this report, 52 sources have been collected and triaged into a source
register (bibliography entry, summary, and 0-5 relevance rating for each);
roughly a quarter have received a full field-by-field deep annotation, with
the remainder prioritized for the literature-review milestone that follows
this interim report. What follows synthesizes the current state of four
literature tracks: (1) benchmark documents describing how NAS-scale
architecture is actually planned today, (2) airline planning and scheduling,
(3) day-of-operations and turnaround coordination, and (4) operations
control, dispatch, and disruption management -- closing with the gap in that
literature this project addresses.

\subsubsection*{Benchmark: how NAS-scale architecture is planned today}

The closest thing to a "best-in-class" comparison for this project's own
deliverable -- a NAS reference architecture -- is the pair of national/
regional ATM modernization roadmaps the FAA and its European counterpart
publish. The FAA's \emph{NAS Enterprise Architecture: Infrastructure
Roadmaps} lays out a 17-year, domain-by-domain (Aircraft, Airport, Airspace
\& Procedures, Automation, Communication, Navigation, Safety, Surveillance,
Weather, and others) investment plan from an "As-Is" to a "To-Be"
NAS~\textcite{faaNasInfrastructureRoadmaps2025}. Its European
counterpart, the SESAR Joint Undertaking's \emph{European ATM Master Plan},
sets out a comparable "Digital European Sky" vision built around
trajectory-based operations (TBO), a service-oriented delivery model, and an
explicit target enterprise architecture
diagram~\textcite{sesarju2025masterPlan}. Both are authoritative and current,
but neither publishes the MBSE derivation behind its architecture -- they
present the resulting roadmap, not a traceable model connecting operational
concept to system architecture.

That derivation step does have a documented precedent, just not from the
NAS: the EU-China GreAT (Greener Air Traffic Operations) research project
publishes a three-volume progression that starts from a current-state TBO
concept baseline~\textcite{great2020d21tboConcept}, uses an explicit MBSE
methodology to derive a gate-to-gate operational and system architecture from
that concept~\textcite{great2021d22operationalArch}, and extends the result
down to an avionics-level architecture. This is a rare, concrete example of
exactly the operational-concept-to-architecture move this capstone's own
methodology (SysML in Cameo, informed by MagicGrid and selected UAF/DoDAF
concepts -- decision D-003) needs to make for the NAS, and the closest
available benchmark for judging whether this project's own architecture is
comparably rigorous. A second, smaller-scale precedent for the same
ontology-first, MBSE-second move applied to an aviation System of Systems is
\textcite{sinharoy2024ontologyUAM}'s architecture framework for Urban Air
Mobility, and \textcite{liu2025mbseAtmSmt} demonstrates a full,
requirements-traceable MBSE pipeline (metamodel, architecture model library,
and formal verification against ICAO requirements) for an ATM system in a
different regulatory context -- together these establish that MBSE-derived,
requirements-traceable ATM architecture is an active, current methodology,
not a hypothetical one, even though no published example does it for the
U.S. NAS specifically.

\subsubsection*{Airline planning and scheduling}

The airline-planning literature converges on one finding relevant to this
project regardless of which sub-problem is being studied:
\textcite{eltoukhy2017airline}'s review of over 100 papers organizes airline
schedule planning into four sequential sub-problems -- flight scheduling,
fleet assignment, aircraft-maintenance routing, and crew scheduling -- solved
one after another, and finds that this sequential decomposition is the
industry norm even though it is provably suboptimal relative to integrated
formulations. \textcite{lohatepanont2004airline} demonstrate a concrete
integrated schedule-design-and-fleet-assignment model, using real airline
data, that outperforms the sequential approach precisely because it does not
treat fleet assignment as a fixed input to be optimized around after schedule
design is already frozen. \textcite{xu2024airlineSchedOpt} extends this
picture with explicit mathematical formulations for thirteen representative
scheduling models and highlights robust-scheduling variants that explicitly
price in disruption risk -- a step toward internalizing, rather than
externalizing, the very costs that later disruption-management research
(below) has to absorb after the fact. The throughline for this project is
that the sequential, compartmentalized planning process is not a strawman:
it is the documented, still-common industry practice, and its cost is that
decisions made at one stage (e.g., fleet assignment) become fixed constraints
the next stage (e.g., crew scheduling) cannot renegotiate even when it would
be globally beneficial to do so.

\subsubsection*{Day-of-operations and turnaround}

Turnaround is the clearest example in this literature of a process that is
technically local (one aircraft, one gate, one set of ground crews) but
whose information dependencies are not: \textcite{schultz2017turnaround}'s
review of turnaround operations and simulation approaches documents how
airport capacity pressure and schedule disruption both concentrate at the
turnaround, and surveys the agent-based, discrete-event, and
cellular-automaton simulation frameworks used to study it.
EUROCONTROL's current \emph{Specification for Airport Collaborative Decision
Making}~\textcite{eurocontrolACDMSpec} is the operational answer to that
information problem: rather than optimizing turnaround locally, it defines a
shared "milestone approach" -- explicit target, estimated, and actual times
for each turnaround event -- that airports, airlines, ground handlers, and
the Network Manager all update and read from a common information model.
A-CDM is, in effect, an existing, deployed example of the kind of
information-sharing architecture this capstone is trying to generalize
across the rest of the NAS: it does not eliminate each stakeholder's local
objective, but it makes the information those objectives depend on visible
across organizational boundaries instead of trapped inside one of them.

\subsubsection*{Operations control, dispatch, and disruption management}

This is the deepest and most directly relevant of the four tracks, and its
findings are woven through Section~\ref{sec:introduction}'s framing above;
this subsection adds the pieces not already covered there.
\textcite{clarke1998irregular}'s foundational 1998 survey of the Airline
Operations Control Center (AOCC) -- still cited as the origin point of this
literature by later reviews such as \textcite{hassanDisruptionReview} --
establishes the AOCC's basic structure (a central control center supported
by maintenance- and station-level control centers) and its guiding decision
rule: return to the planned schedule as fast as possible, subject to crew and
maintenance feasibility, largely independent of the trajectory-level fuel or
routing consequences of how that return is achieved.
\textcite{hassanDisruptionReview}'s more recent review formalizes why that
rule persists: aircraft, crew, and passenger recovery are treated as
separable decision domains not because they are independent in reality, but
because integrated (multi-resource) recovery is harder to solve and the
literature that could support it does not match AOCC practice --
\citeauthor{hassanDisruptionReview} describe this explicitly as a gap between
published research and the reality of AOCC decision-making.
\textcite{bouarfa2018masDisruption} and its originating
thesis~\textcite{castro2013aoccMasThesis} offer a partial counterexample: the
MASDIMA multi-agent architecture assigns each resource (aircraft, crew,
passengers) its own agent with an explicit utility function and negotiates a
joint solution through a supervisor agent -- a working, implemented
demonstration that local objectives \emph{can} be reconciled algorithmically
rather than left to sequential handoff, though \citeauthor{bouarfa2018masDisruption}'s
own evaluation finds that automating those roles concentrates workload and
judgment onto the remaining human supervisor rather than eliminating the
need for it.

The pattern above is not limited to human decision processes; it recurs even
in mathematically rigorous work. \textcite{santana2023arpReview}'s systematic
review of the Aircraft Recovery Problem finds that studies (and, by
extension, the AOCC processes they model) rarely minimize delay,
cancellations, and aircraft swaps simultaneously -- they treat these as
separable sub-problems solved and handed off in sequence, so a solution that
is optimal at one stage can be infeasible once the next stage's constraints
are applied. \textcite{hu2024disruptionOptReview} adds that the standard
simplification in this literature -- modeling delay cost as linear in delay
minutes -- does not match airline practice, where delay cost is nonlinear
and airlines routinely trade mathematical optimality for decision speed.
\textcite{kontodimou2026turnaroundBuffer}'s predictive turnaround-buffer
model achieves a real, measured 37.1\% reduction in propagated delay, but
only by holding routing, fleet assignment, crew duty, maintenance, and
passenger connectivity fixed as external constraints -- the improvement is
genuine within its bounded scope, and the boundary itself is the risk this
capstone investigates. Finally, \textcite{dispatcherWorkload2025}'s recent
(2026) study of the flight-dispatcher workload-balancing problem is notable
for two reasons: it is one of only three studies the authors could find that
treat dispatcher scheduling as an optimization problem at all, and it
supplies a directly reusable formulation (a load-based objective, a
Lagrangian-relaxation bound, and a metaheuristic) for treating human workload
as a first-class, quantified constraint rather than an unmeasured
externality -- exactly the kind of factor Section~\ref{sec:overview}
identifies as "harder to quantify" and therefore routinely omitted from
today's cost models.

\subsubsection*{Research gap}

Individually, each of these tracks is well studied: sequential-versus-
integrated optimization within airline planning, information-sharing
protocols for turnaround, and disruption-recovery decision support each have
mature, if incomplete, literatures, and precedents exist for deriving
MBSE-based ATM architecture from an operational concept in other regions
(GreAT, SESAR) and other aviation domains (UAM). What none of them do,
individually or together, is connect these tracks into one traceable
model. Each track studies its own slice of the trajectory-intent chain
in isolation -- planning \emph{or} turnaround \emph{or} dispatch \emph{or}
foreign-ATM architecture -- using that slice's own vocabulary and cost model,
and each documents, in its own domain, the same underlying pattern: sequential
or siloed decision-making that is locally rational and globally uncertain.
No source in this register builds a single architecture that traces
operational intent end-to-end -- from an enterprise objective, through
airline planning, day-of-operations control, and air traffic control, to the
trajectory an aircraft actually flies -- while keeping each stakeholder's
objectives, and the points where one stakeholder's decision imposes a cost on
another, explicit and inspectable in the same model. That is the gap this
capstone's reference architecture is intended to fill: not a new
optimization method for any one of these tracks, but a structure that makes
the couplings \emph{between} them visible to a systems engineer or
decision-support tool before a disruption forces them into view.


\subsection{Stakeholders}
\label{sec:stakeholders}

Stakeholders were identified with a PESTLE (Political, Economic, Social,
Technical, Legal, Environmental) sweep across the NAS, producing a register
of several dozen stakeholder classes ranging from the FAA and ICAO down to
local communities and future generations affected by long-run environmental
cost. That full inventory is deliberately broader than any one model or
optimization study can carry in detail, so it is narrowed, for the
architecture and objective-ontology work, to seven categories that sit
directly inside the trajectory-intent chain and actively trade against one
another: \textbf{Airline}, \textbf{Passenger}, \textbf{ATC/ANSP},
\textbf{Airport}, \textbf{Flight Crew}, \textbf{Environmental/Societal}, and
\textbf{Military}. Each of the seven has been carried through a per-category
objective/cost/value analysis (classifying each objective as a hard
constraint, an optimization objective, a cost/penalty, a measure of
performance, or a measure of effectiveness, and rating whether it is
directly, indirectly, or not coupled to aircraft trajectory) that is
summarized here and used directly in Section~\ref{sec:scope} and
Section~\ref{sec:deliverables}:

\begin{itemize}
  \item \textbf{Airline} -- fuel cost, crew cost, aircraft utilization,
    maintenance, schedule integrity, passenger connections, delay cost,
    revenue, and dispatch reliability. Represented operationally by the
    AOCC, dispatch, and network-planning functions discussed in
    Section~\ref{sec:litreview}.
  \item \textbf{Passenger} -- ticket price, travel time, connection
    reliability, schedule convenience, comfort, and disruption risk. A
    comparatively low-power stakeholder in the sense that passengers do not
    make operational decisions, but one whose travel-time and connection
    outcomes are a direct, measurable consequence of decisions made by every
    other stakeholder in this list.
  \item \textbf{ATC/ANSP} -- safety and separation, sector workload,
    capacity, predictability, delay, and traffic complexity. Safety and
    separation are hard constraints that bound every other objective in this
    architecture; sector workload is the clearest example in this project of
    an externality (imposed by airline and individual-aircraft trajectory
    requests) that the imposing party does not itself pay for.
  \item \textbf{Airport} -- gate utilization, runway utilization, surface
    congestion, turnaround performance, and passenger throughput. Sits at
    the A-CDM information-sharing boundary discussed in
    Section~\ref{sec:litreview}.
  \item \textbf{Flight Crew} -- safety, workload, procedural compliance,
    schedule/duty constraints, and operational flexibility. The Captain's
    authority to deviate from an assigned trajectory to preserve safety
    margin is this project's clearest example of decision authority that
    can, and by design does, override every other stakeholder's local
    objective.
  \item \textbf{Environmental/Societal} -- CO2, NOx, contrail/climate
    effects, noise, local air quality, and community impacts. The
    lowest-power stakeholder category in this list -- represented by
    advocacy groups, regulators, and "future generations" as a societal
    proxy rather than by an operational decision-maker -- and the clearest
    example of a stakeholder whose costs are chronically externalized by
    every other category's schedule-recovery and efficiency decisions.
  \item \textbf{Military} -- mission effectiveness, mission timing, fuel
    availability, airspace access, operational security, and resilience.
    Included specifically at the civil-military airspace-access interface
    (special-use-airspace activation forcing civil rerouting), not as a
    detailed model of military mission planning -- see
    Section~\ref{sec:scope}.
\end{itemize}

Because this is a single-student capstone rather than a team project, there
is no team-composition bias to check against; the mitigation instead was to
seed the register broadly (the full PESTLE inventory) before narrowing to
the seven working categories, so that low-power stakeholders were considered
and then deliberately retained -- Passenger and Environmental/Societal in
particular -- rather than dropped for convenience. The needs, resources,
constraints, and information requirements underlying each category above
are being developed into a full RACCI (Responsible, Accountable, Consulted,
Informed, and a fifth authority dimension) matrix of NAS operational
decisions; that matrix, and the deeper per-stakeholder detail, are in
progress and will be reported in a subsequent submission.

\subsection{Scope}
\label{sec:scope}

The System of Interest (SOI) is the portion of the NAS that participates in
the trajectory-intent chain for scheduled U.S. commercial (Part 121) air
transportation: the organizations, systems, and information exchanges that
carry intent from an airline's strategic and network-planning objectives,
through day-of-operations control and air traffic control, to the trajectory
an aircraft actually flies and the guidance commands that execute it.

\textbf{Included.} Airline strategic/network planning and its
sequentially-coupled sub-problems (schedule design, fleet assignment,
aircraft routing, crew scheduling); day-of-operations functions (dispatch,
the AOCC, weight and balance); turnaround and airport ground operations
insofar as they exchange information across organizational boundaries (the
A-CDM milestone model); air traffic control from clearance delivery through
enroute, arrival, and approach control; and the onboard systems (FMS,
autopilot, guidance) that translate a cleared trajectory into aircraft
motion. The seven stakeholder categories in Section~\ref{sec:stakeholders}
are the SOI's actor set.

\textbf{Excluded.} General aviation, all-cargo and charter operations
outside scheduled Part 121 service, and international airspace beyond the
NAS boundary are out of scope for direct modeling, though the SESAR/European
and GreAT sources in Section~\ref{sec:litreview} are used as comparative
benchmarks, not as SOI content. Military aviation is excluded except at its
one interface with the civil SOI: special-use-airspace activation and the
civil rerouting it forces (Section~\ref{sec:stakeholders}); internal military
mission planning, weapons systems, and mission systems are treated as an
opaque boundary condition. Urban Air Mobility and uncrewed systems are
excluded from the SOI itself and used only as methodological precedent
(Section~\ref{sec:litreview}). Below the aircraft/FMS interface --
autopilot control-law design, control-surface actuation, and other
airframe-internal behavior -- is excluded; the architecture treats "aircraft
motion" as the observable output of guidance commands, not as a flight-
dynamics model.

\textbf{Level of abstraction.} The model spans enterprise objectives down to
trajectory intent and the guidance commands that realize it, but not below
that: the trajectory-intent chain (Strategic objective $\rightarrow$ mission
plan $\rightarrow$ flight plan $\rightarrow$ ATC constraints $\rightarrow$
trajectory negotiation $\rightarrow$ FMS intent $\rightarrow$ guidance
commands $\rightarrow$ aircraft motion, Section~\ref{sec:overview}) is
modeled at full fidelity end-to-end, while lower-level detail inside any one
link -- e.g., the specific numerical control law an FMS uses to track a
cleared trajectory -- is treated as a boundary input. Consistent with the
per-objective abstraction judgments already made in the stakeholder
objective/cost ontology (Section~\ref{sec:stakeholders}), an objective or
decision is modeled at full trajectory fidelity when it is \emph{directly}
coupled to trajectory (e.g., ATC safety/separation, schedule-recovery delay
cost, contrail avoidance) and treated as a coarser boundary condition when
the coupling is indirect (e.g., ticket pricing, passenger throughput, crew
duty-time constraints).

\textbf{Boundary criteria.} A candidate element is included as a
\emph{system} in the architecture if it produces, transforms, transmits, or
consumes trajectory-relevant intent or information, and if a specific
organization can be identified as holding authority or responsibility for it
-- consistent with the domain decomposition already adopted (governance,
airspace management, airspace resources, flight operations, airport
operations, aircraft systems, information services, infrastructure, and
decision support; decision D-002). An element that merely exists in the
aviation domain without a traceable connection to that chain (e.g., an
airline's internal HR system) is out of scope. A candidate \emph{stakeholder}
or actor is included if it holds authority over, is responsible for
executing, or bears a measurable cost or benefit from a decision inside the
trajectory-intent chain -- the test applied when the seven-category working
set in Section~\ref{sec:stakeholders} was narrowed from the full PESTLE
inventory.

\textbf{Status.} This is the project's initial, working SOI boundary,
consistent with the candidate scope already sketched in the project's
working notes; it closes the boundary-definition questions raised during
research-framework setup (inclusion/exclusion, abstraction level, and system
and stakeholder boundary criteria). It remains open to revision when
alternative decompositions are compared against it later in the project
(organization-based, lifecycle-based, physical, information-flow, and
decision-authority views), and that comparison may refine or replace the
domain decomposition without changing the SOI boundary stated here.

\subsection{Project Deliverables}
\label{sec:deliverables}

The project has two primary deliverables, developed together so that the
second demonstrates the value of the first.

\textbf{1. A SysML/Cameo reference architecture of the SOI defined in
Section~\ref{sec:scope}.} This includes a NAS context diagram and
stakeholder model; block definition diagrams (BDDs) for the domain
decomposition; internal block diagrams (IBDs) for the critical interfaces
identified along the trajectory-intent chain; an information-object and
item-flow model; a nominal-flight activity diagram and responsibility
swimlanes (grounded in the ConOps built from Section~\ref{sec:litreview}'s
literature, especially the A-CDM milestone model and the AOCC/dispatch
decision process); and a requirements model traced from stakeholder needs
(Section~\ref{sec:stakeholders}) through requirements, systems, activities,
and information exchanges, to the decisions that ultimately shape aircraft
behavior.

\textbf{2. A bounded optimization-demonstration capability} built as one
decision-support capability \emph{inside} that architecture rather than as
a standalone study -- consistent with this project's own framing (this is a
systems-engineering capstone, not an optimization paper). Using the
objective/cost ontology summarized in Section~\ref{sec:stakeholders}, the
demonstration formulates one representative multi-stakeholder conflict (the
leading candidate is airline fuel-cost minimization versus ATC/ANSP sector
workload, identified in that ontology as the cleanest available decision
variable with clean opposing measures on both sides), runs a single
local-versus-system-level-optimum comparison and objective-weight sweep, and
traces the resulting decision back through the architecture to show which
stakeholders absorb which costs -- the specific capability a standalone
optimizer, decoupled from the architecture, could not show.

Success is assessed against measures that can be checked directly in the
finished model and report, not by subjective judgment of "improvement":

\begin{itemize}
  \item \textbf{Stakeholder coverage} -- all seven stakeholder categories
    (Section~\ref{sec:stakeholders}) represented in the architecture with an
    identified objective, decision authority, and at least one information
    dependency.
  \item \textbf{Traceability coverage} -- a complete, unbroken chain from
    stakeholder need through requirement, system, activity, and information
    exchange, to a specific effect on aircraft trajectory or behavior, for
    at least the nominal-flight ConOps scenario.
  \item \textbf{Interface/requirement completeness} -- all interfaces
    identified as critical in the IBD model are connected, and all
    requirements trace to at least one stakeholder need.
  \item \textbf{Demonstrated tradeoff} -- the optimization capability
    produces at least one local-optimum-versus-SoS-optimum comparison that
    quantifies a cost externalized from one stakeholder onto another, and
    shows that externality changing as the objective weighting changes.
  \item \textbf{Architecture value-add} -- an explicit statement, supported
    by the demonstration, of what the architecture exposes that the
    optimization result alone would not (i.e., which stakeholder absorbs the
    externality and through which interface).
\end{itemize}

\subsection{Project Timing}
\label{sec:timing}

The timeline below revisits the milestone schedule from the 503 Proposal
submittal and maps each milestone to the corresponding section(s) of the
project's working to-do list (16 sections, research framework through final
deliverables). Three planning activities (alternative decomposition
exploration, myopic-optimization conflict analysis, and the capstone
demonstration walkthrough) have no dedicated milestone in the original
submittal; they are scheduled to fall between their logically adjacent
milestones and are marked accordingly.

\begin{longtable}{p{2.3cm}p{6.3cm}p{1.6cm}p{3.7cm}}
\toprule
\textbf{Target date} & \textbf{Milestone} & \textbf{Plan \S} & \textbf{Status (as of \today)} \\
\midrule
\endhead
2026-09-07 & PESTLE Analysis \& System Definition (stakeholder discovery, SOI boundary, NAS context diagram) & S1, S7, S10 & PESTLE discovery and stakeholder register complete; SOI boundary stated in Section~\ref{sec:scope} \\
2026-09-10 & Authority \& Responsibility Models (RACCI, responsibility swimlanes) & S7, S10 & Not started \\
\textbf{2026-09-20} & \textbf{Interim Report \#1 (this report)} & S1--S4, S16 & In progress -- this document \\
2026-09-22 & Sequence Diagrams (critical message flows) & S10 & Not started \\
2026-10-01 & Activity Diagrams (nominal-flight scenario walkthrough) & S10 & Not started \\
2026-10-12 & Requirement Traceability (needs $\to$ requirements $\to$ systems $\to$ activities $\to$ information) & S10 & Not started \\
2026-10-17 & Measures of Effectiveness (MOE/MOP hierarchy) & S8 & Complete -- see stakeholder objective/cost ontology \\
2026-10-20 & Requirements Definition & S10 & Not started \\
(unscheduled, between 2026-10-17 and 2026-11-03) & Alternative system decompositions compared against the domain decomposition & S6 & Not started \\
2026-10-25 & Internal Block Diagrams (critical interfaces connected) & S10 & Not started \\
2026-10-29 & Interface Definition (Interface Control Document) & S10 & Not started \\
2026-11-03 & Draft Architecture (BDDs \& package structure) & S10 & Not started \\
2026-11-08 & Operational Scenarios (nominal-flight ConOps document) & S5 & Not started \\
(unscheduled, between MOE and Optimization Prototype) & Myopic-optimization conflict analysis & S9 & Not started \\
2026-11-17 & Literature Review (annotated bibliography, 10--20 papers categorized) & S2--S4 & In progress -- 52 sources triaged, deep annotation ongoing \\
2026-11-21 & Optimization Prototype (working code, unit tests) & S11, S12 & Not started \\
2026-11-24 & Integration of MBSE Architecture (traceability links, end-to-end scenario) & S10, S13 & Not started \\
(unscheduled, between Integration and Final Paper) & Capstone demonstration walkthrough & S14 & Not started \\
2026-12-18 & Final Paper (complete draft) & S16 & Not started \\
2026-12-20 & Paper Submittal (official) & S16 & Not started \\
2026-12-26 & Final Baseline \& Draft Paper (versioned model \& draft paper) & S15 & Not started \\
\bottomrule
\end{longtable}

Two scheduling notes carry over from the original submittal and are flagged
here rather than silently resolved: the submittal schedules the Literature
Review milestone (2026-11-17) \emph{after} most of the architecture
milestones (2026-09-07 through 2026-11-03), later than this report's own
literature-review-before-architecture sequencing assumes -- in practice,
literature work is running continuously in the background (as reflected in
Section~\ref{sec:litreview}) rather than strictly front-loaded. The Final
Baseline \& Draft Paper milestone (2026-12-26) is also dated \emph{after}
Paper Submittal (2026-12-20); this is being treated as an intentional
post-submission archival/versioning step pending confirmation with the
faculty adviser.
